\chapter{Bloch Geometry, Euler Phase, and Spinorial Closure}
\label{chap:spinor}

\section{Bloch-sphere parameterization}

Every normalized pure state in $\C^2$, up to global phase, can be written as
\begin{equation}
 \ket{\psi(\theta,\phi)}
 =\cos\frac{\theta}{2}\ket{0}
 +e^{i\phi}\sin\frac{\theta}{2}\ket{1},
 \qquad 0\le\theta\le\pi.
 \label{eq:bloch-state}
\end{equation}
Comparing with \cref{eq:binary-state},
\begin{equation}
 \sigma=\cos^2\frac{\theta}{2},
 \qquad 1-\sigma=\sin^2\frac{\theta}{2}.
 \label{eq:sigma-theta}
\end{equation}
Therefore
\begin{equation}
 \sigma=\frac12\quad\Longleftrightarrow\quad\theta=\frac\pi2.
 \label{eq:half-equator}
\end{equation}
Balanced binary states form the Bloch equator. Exact equal-gain cancellation adds $\phi=\pi$, selecting the ray $\ket{-}$.

\begin{figure}[htbp]
 \centering
 \includegraphics[width=0.70\textwidth]{bloch_equator.png}
 \caption{The half condition fixes the equator, while the relative phase $\pi$ selects the state orthogonal to the equal-gain detector.}
 \label{fig:bloch-equator}
\end{figure}

\section{Euler's sign}

Euler's identity
\[
 e^{i\pi}=-1
\]
is the elementary phase statement behind destructive interference. If a normalized parameter $\sigma\in(0,1)$ controls a phase fraction through
\begin{equation}
 U(\sigma)=e^{2\pi i\sigma},
 \label{eq:phase-law}
\end{equation}
then
\begin{equation}
 U(\sigma)=-1
 \quad\Longleftrightarrow\quad
 \sigma=\frac12
 \quad\text{for }0<\sigma<1.
 \label{eq:phase-half}
\end{equation}
This is an exact implication after the phase law is defined. The non-trivial assumption is the identification of the same $\sigma$ with the normalized phase fraction and, later, with $\Re s$.

\section{Spinorial sign under $2\pi$ rotation}

The group $\mathrm{SU}(2)$ double-covers $\mathrm{SO}(3)$. For a spin-$1/2$ representation, a physical rotation through angle $\alpha$ about unit vector $\mathbf n$ is represented by
\begin{equation}
 U(\alpha,\mathbf n)=\exp\!\left(-\frac{i\alpha}{2}\,\mathbf n\cdot\boldsymbol\sigma\right),
 \label{eq:spin-rotation}
\end{equation}
where $\boldsymbol\sigma$ are the Pauli matrices. At $\alpha=2\pi$,
\begin{equation}
 U(2\pi,\mathbf n)=-I,
 \label{eq:spinor-minus}
\end{equation}
and only after $4\pi$ does the state vector return to itself. The global sign is invisible for an isolated ray but becomes observable in interference.

Berry's geometric phase generalizes this phenomenon. For a spin-$1/2$ eigenstate transported adiabatically around a closed loop on the Bloch sphere, the geometric phase is minus one half the enclosed solid angle, modulo $2\pi$ \citep{berry1984}. A loop enclosing solid angle $2\pi$ therefore produces a phase $\pi$ and the sign $-1$.

The relevance to the present ansatz is structural: the same relative sign required for exact two-channel cancellation is the sign associated with spinorial closure. This does not show that zeta zeros are physical spinors. It identifies a two-dimensional projective geometry in which balance and sign closure meet naturally.

\section{The binary-spinor equivalence theorem}

Define the phase-locked amplitude
\begin{equation}
 \Acomp_*(\sigma)
 =\sqrt\sigma+e^{2\pi i\sigma}\sqrt{1-\sigma},
 \qquad 0<\sigma<1.
 \label{eq:phase-locked-amplitude}
\end{equation}

\begin{theorem}[Triple coincidence of the half]
\label{thm:triple-coincidence}
For $0<\sigma<1$, the following statements are equivalent:
\begin{enumerate}
  \item $\sigma=1-\sigma$;
  \item $\Hbin(\sigma)=\ln2$;
  \item $e^{2\pi i\sigma}=-1$;
  \item $\Acomp_*(\sigma)=0$;
  \item $\sigma=1/2$.
\end{enumerate}
\end{theorem}
\begin{proof}
The complement fixed-point equation gives $\sigma=1/2$. By \cref{lem:entropy-max}, the entropy reaches $\ln2$ only there. Equation \cref{eq:phase-half} establishes the phase equivalence. Finally, \cref{lem:exact-cancellation} applied to $\phi=2\pi\sigma$ gives the amplitude equivalence.
\end{proof}

This theorem is one of the main exact results of the monograph. It shows that four apparently different notions collapse onto the same point inside a single explicit model. The theorem says nothing yet about $\zeta$ or $\xi$; it becomes relevant to RH only if $\sigma$ is canonically tied to $\Re s$ and zerohood is tied to the readout.

\section{Gauge and global phase}

A pure quantum state is a ray, so multiplication by $e^{i\chi}$ does not change the state. The cancellation amplitude, however, compares a state with a fixed detector and is sensitive to relative phase. Therefore a candidate zeta-state map must specify enough gauge structure to make the readout meaningful.

Suppose
\[
 \ket{\Psi(s)}\mapsto e^{i\chi(s)}\ket{\Psi(s)}.
\]
Then $\braket{+}{\Psi(s)}$ acquires the same factor, so its zero set is gauge invariant. This is favourable: an amplitude-zero condition can be well-defined on rays. But decomposing the state into two channel phases may depend on basis and gauge. The basis must therefore arise from the complement symmetry rather than arbitrary coordinates.

\section{Covariance and the equator}

The channel-swap covariance previewed in \cref{eq:state-covariance-preview} exchanges the north and south hemispheres under $\sigma\mapsto1-\sigma$. On the fixed line it places the state on the equator, because the two component norms must agree. A separate analytic or spectral condition must select the antipodal phase required for orthogonality to $\ket{+}$.

This separation is useful:
\begin{align*}
 \text{zeta involution covariance}&\quad\rightsquigarrow\quad\text{equal component norms on }\Crit,\\
 \text{zero condition}&\quad\rightsquigarrow\quad\text{relative phase }\pi,\\
 \text{together}&\quad\rightsquigarrow\quad\text{exact cancellation}.
\end{align*}
A successful bridge should derive both arrows from one construction.

\begin{warning}
The spinorial sign is not a proof of the phase law $\phi=2\pi\Re s$. That law is a model identification. Any future derivation must explain why the zeta argument traverses a loop whose holonomy is parameterized by the real part in precisely this way.
\end{warning}

\status{Bloch parameterization, spinor sign, Berry-phase relation, and the triple-coincidence theorem are exact within their stated model. The identification with zeta data remains open.}
